\documentclass{beamer}
\usepackage{graphicx}
\usepackage{xcolor}
\usepackage{tikz}

% 自定义颜色
\definecolor{purpleblue}{RGB}{80, 80, 180}
\definecolor{lightgray}{RGB}{230, 230, 230}
\definecolor{novelPurple}{RGB}{98, 54, 255} 
\definecolor{darkgray}{RGB}{50, 50, 50}

% 主题设置
\usetheme{Madrid}  
\usecolortheme{dolphin} 

% 设置字体（移除 fontspec，确保 pdfLaTeX 兼容）
\setbeamerfont{title}{size=\Large,series=\bfseries}
\setbeamerfont{frametitle}{size=\LARGE,series=\bfseries}
\setbeamercolor{frametitle}{fg=white, bg=novelPurple}
\setbeamercolor{title}{fg=white, bg=purpleblue}

% 封面信息
\title[SAT Hard Question]{\textbf{SAT Hard Question 4}}
\author{\textbf{Jack Zeng}}
\institute{\textcolor{white}{\textbf{Novel Prep}}}
\date{\today} 

\begin{document}

% --- Title Slide ---
\begin{frame}
    \titlepage
    \vfill
    \begin{flushright}
        \textcolor{purpleblue}{\Huge \textbf{Novel Prep}}
    \end{flushright}
\end{frame}


\begin{frame}
    \begin{tikzpicture}[remember picture, overlay]
        % 设置浅色渐变背景
        \shade[top color=softPurple, bottom color=softGray] 
            (current page.north west) rectangle (current page.south east);
        
        % 添加高斯模糊光晕
        \fill[white, opacity=0.3] (current page.center) circle (4cm);
        
        % 主标题
        \node[align=center] at (current page.center) {
            {\Huge\bfseries\textcolor{purpleblue}{Novel Prep}} \\[0.5cm]
            {\Large\itshape\textcolor{lightText}{Excellence in SAT Prep}}
        };

        % 添加柔和光点（点缀）
        \foreach \x in {1,2,3,4,5} {
            \fill[purpleblue, opacity=0.2] (rand*10-5, rand*7-3.5) circle (0.1+\x*0.05);
        }

        ;
    \end{tikzpicture}
\end{frame}

\begin{frame}{\textbf{Question 1: Probability (Statistics)}}
\small
The incomplete table below summarizes the number of left-handed and right-handed students by gender for eighth-grade students at Keisel Middle School.

\begin{center}
\begin{tabular}{|c|c|c|}
\hline
\textbf{Gender} & \textbf{Left} & \textbf{Right} \\
\hline
Female &  &  \\
\hline
Male &  &  \\
\hline
Total & 18 & 122 \\
\hline
\end{tabular}
\end{center}

There are \textbf{5 times} as many right-handed female students as left-handed female students, and \textbf{9 times} as many right-handed male students as left-handed male students.

If there are a total of 18 left-handed and 122 right-handed students, which of the following is closest to the probability that a randomly selected right-handed student is female?

\begin{itemize}
    \item[A.] \( 0.410 \)
    \item[B.] \( 0.357 \)
    \item[C.] \( 0.333 \)
    \item[D.] \( 0.250 \)
\end{itemize}
\end{frame}

\begin{frame}{\textbf{Answer to Question 1}}
Let \( x \) be the number of left-handed females and \( y \) be the number of left-handed males.

\begin{itemize}
    \item Total left-handed: \( x + y = 18 \)
    \item Right-handed females: \( 5x \), right-handed males: \( 9y \)
    \item Total right-handed: \( 5x + 9y = 122 \)
\end{itemize}

Now solve the system:
\[
\begin{aligned}
x + y &= 18 \\
5x + 9y &= 122
\end{aligned}
\]

Substitute \( y = 18 - x \) into the second equation:
\[
5x + 9(18 - x) = 122 \Rightarrow 5x + 162 - 9x = 122 \Rightarrow -4x = -40 \Rightarrow x = 10
\]
\[
\Rightarrow y = 8, \quad \text{Right-handed females} = 50, \quad \text{Right-handed students} = 122
\]

\[
\text{Probability} = \frac{50}{122} \approx 0.41
\]

\textbf{Correct Answer: A. \( 0.410 \)}
\end{frame}



\begin{frame}{\textbf{Question 2: Little's Law}}
If shoppers enter a store at an average rate of \( r \) shoppers per minute and each stays in the store for an average of \( T \) minutes, the average number of shoppers in the store is \( N = rT \).

In the original store:
\[
r = 3 \text{ shoppers/minute}, \quad T = 15 \Rightarrow N = 45
\]

In the new store:
\[
r = \frac{90}{60} = 1.5, \quad T = 12 \Rightarrow N = 1.5 \times 12 = 18
\]

What percent less is the new store’s average number of shoppers compared to the original?

\begin{itemize}
    \item[A.] \( 57 \)
    \item[B.] \( 58 \)
    \item[C.] \( 59 \)
    \item[D.] \( 60 \)
\end{itemize}
\end{frame}

\begin{frame}{\textbf{Answer to Question 2}}
From Little's Law:
\[
\text{Original store: } N = 3 \times 15 = 45
\]
\[
\text{New store: } r = \frac{90}{60} = 1.5,\quad T = 12,\quad N = 1.5 \times 12 = 18
\]

\[
\text{Percent decrease} = \frac{45 - 18}{45} \times 100 = \frac{27}{45} \times 100 = 60\%
\]

\textbf{Correct Answer: D. \( 60 \)}
\end{frame}

\begin{frame}{\textbf{Question 3: Statistical Generalization}}
\small
A trivia tournament organizer wanted to study the relationship between the number of points a team scores in a trivia round and the number of hours that a team practices each week.

For the study, the organizer selected \textbf{55} teams at random from all trivia teams in a certain tournament. The table displays the information for the \textbf{40} teams in the sample that practiced for at least \textbf{3 hours per week.}

\begin{center}
\begin{tabular}{|c|c|c|c|}
\hline
\textbf{Hours practiced} & \textbf{6 to 13 points} & \textbf{14 or more points} & \textbf{Total} \\
\hline
\textbf{3 to 5 hours} & 6 & 4 & 10 \\
\textbf{More than 5 hours} & 4 & 26 & 30 \\
\hline
\textbf{Total} & 10 & 30 & 40 \\
\hline
\end{tabular}
\end{center}

Which of the following is the \textbf{largest population} to which the results of the study can be generalized?

\begin{itemize}
    \item[A.] All trivia teams in the tournament that scored \textbf{14 or more points}.
    \item[B.] The \textbf{55 trivia teams} in the sample.
    \item[C.] The \textbf{40 trivia teams} in the sample that practiced for at least \textbf{3 hours per week}.
    \item[D.] All trivia teams in the \textbf{tournament}.
\end{itemize}
\end{frame}

\begin{frame}{\textbf{Answer to Question 3}}


\textbf{Correct Answer: D. }
\end{frame}

\begin{frame}{\textbf{Question 4: Tangent Line and Circle Geometry}}
\textbf{Given:} The circle has center \( (-1, 1) \). Line \( t \) is tangent to this circle at point \( (4, -3) \).

Which of the following points lies on line \( t \)?

\begin{itemize}
    \item[A.] \( (0, \frac{5}{4}) \)
    \item[B.] \( (3,6) \)
    \item[C.] \( (8,2) \)
    \item[D.] \( (9,1) \)
\end{itemize}
\end{frame}

\begin{frame}{\textbf{Answer to Question 4}}
To find the tangent line to a circle at a point, find the radius vector from the center to the point of tangency:
\[
\vec{r} = \langle 4 - (-1), -3 - 1 \rangle = \langle 5, -4 \rangle
\]
The slope of the radius is \( \frac{-4}{5} \), so the slope of the tangent line is the negative reciprocal: \( \frac{5}{4} \).

Using point-slope form:
\[
y + 3 = \frac{5}{4}(x - 4)
\]

Now test which point lies on this line. Try \( (8, 2) \):
\[
2 + 3 = \frac{5}{4}(8 - 4) \Rightarrow 5 = \frac{5}{4} \cdot 4 = 5
\]

\textbf{Correct Answer: C. \( (8,2) \)}
\end{frame}

\begin{frame}{\textbf{Question 5: Rectangular Prism Surface Area}}
Two identical rectangular prisms each have a height of \textbf{90 cm}. The base of each prism is a square, and the surface area of each prism is \( K \) cm\(^2\).

When the prisms are glued together along one square base, the new surface area is:
\[
\frac{92}{47} K \text{ cm}^2.
\]

What is the side length, in cm, of each square base?

\begin{itemize}
    \item[A.] \( s = 4 \)
    \item[B.] \( s = 6 \)
    \item[C.] \( s = 8 \)
    \item[D.] \( s = 10 \)
\end{itemize}
\end{frame}


\begin{frame}{\textbf{Answer to Question 5}}
Let the side length be \( s \). For a rectangular prism with square base:
\[
\text{Surface Area} = 2s^2 + 4sh
\]

Each original prism has:
\[
K = 2s^2 + 4s(90) = 2s^2 + 360s
\]

When glued together along one square base, one face is removed:
\[
\text{New surface area} = 2K - 2s^2 = 2(2s^2 + 360s) - 2s^2 = 2s^2 + 720s
\]

Given:
\[
2s^2 + 720s = \frac{92}{47}(2s^2 + 360s)
\]

Solve for \( s \), and you’ll find \( s = 8 \)

\textbf{Correct Answer: C. \( s = 8 \)}
\end{frame}






\begin{frame}{\textbf{Question 6: Triangle Similarity and Area}}
In triangle \( RST \), angle \( T \) is a right angle. Point \( L \) lies on \( \overline{RS} \), point \( K \) lies on \( \overline{ST} \), and \( \overline{LK} \) is parallel to \( \overline{RT} \).

\textbf{Given:}
\begin{itemize}
    \item \( RT = 72 \) units,
    \item \( LK = 24 \) units,
    \item Area of \( \triangle RST = 792 \) square units.
\end{itemize}

Find the length of \( KT \), in units.
\end{frame}


\begin{frame}{\textbf{Answer to Question 6}}
Since \( LK \parallel RT \), triangle \( \triangle LKT \sim \triangle RST \) by AA similarity.

\[
\frac{LK}{RT} = \frac{24}{72} = \frac{1}{3}
\]

Because area scales by the square of the side ratio:
\[
\text{Area of } \triangle LKT = \left( \frac{1}{3} \right)^2 \cdot 792 = \frac{1}{9} \cdot 792 = 88
\]

Now use triangle area formula:
\[
\text{Area} = \frac{1}{2} \cdot LK \cdot KT = 88
\Rightarrow \frac{1}{2} \cdot 24 \cdot KT = 88
\Rightarrow 12 \cdot KT = 88
\Rightarrow KT = \frac{88}{12} = \frac{22}{3}
\]

\textbf{Correct Answer: \( \boxed{\frac{22}{3}} \) units}
\end{frame}



\begin{frame}{\textbf{Question 7: Circle Equation (Geometry)}}
Given the circle equation:
\[
(x - 6)^2 + (y + 5)^2 = 16
\]

Point \( P(10, -5) \) lies on the circle, and segment \( PQ \) is the diameter.

\begin{itemize}
    \item What are the coordinates of point \( Q \)?
\end{itemize}
\end{frame}

\begin{frame}{\textbf{Answer to Question 7}}
\begin{itemize}
    \item The center of the circle is the midpoint of diameter \( PQ \), and is given as:
    \[
    (6, -5)
    \]
    
    \item Let point \( Q = (x, y) \), and recall that the midpoint of \( PQ \) is:
    \[
    \left( \frac{10 + x}{2}, \frac{-5 + y}{2} \right) = (6, -5)
    \]
    
    \item Solve for \( x \) and \( y \):
    \[
    \frac{10 + x}{2} = 6 \Rightarrow 10 + x = 12 \Rightarrow x = 2
    \]
    \[
    \frac{-5 + y}{2} = -5 \Rightarrow -5 + y = -10 \Rightarrow y = -5
    \]
\end{itemize}

\textbf{Coordinates of point \( Q \):} \( \boxed{(2, -5)} \)
\end{frame}




\end{document}